% ID: SEL-D0022-OB-E01
% ESTADO: validada. S023, revisada con los criterios oficiales.
Sean $x,y$ los litros de colonia A y B. Las restricciones de extracto
y alcohol son $0.05x+0.10y\le70$ y $0.10x+0.15y\le120$.
Equivalentemente, $x+2y\le1400$ y $2x+3y\le2400$, con $x,y\ge0$.
En los ejes quedan $(1200,0)$ y $(0,700)$. Al resolver las dos
ecuaciones frontera simultáneamente se obtiene $(600,400)$.
\begin{center}
\setlength{\unitlength}{5mm}
\begin{picture}(16,10)(-2,-2)
\qbezier(0,0)(7,0)(14,0)\put(14.2,0){$x$}
\qbezier(0,0)(0,3.8)(0,7.6)\put(0,8){$y$}
\qbezier(0,7)(3,5.5)(6,4)
\qbezier(6,4)(9,2)(12,0)
\put(0,6.6){\makebox(0,0)[r]{$(0,700)$}}
\put(6,4.4){$(600,400)$}
\put(12,-1){\makebox(0,0){$(1200,0)$}}
\end{picture}
\end{center}
El ingreso $24x+40y$ es lineal. En $(0,0)$, $(1200,0)$, $(600,400)$
y $(0,700)$ vale, respectivamente, 0, 28800, 30400 y 28000 euros.
El máximo se obtiene al fabricar
$\boxed{600\text{ L de A y }400\text{ L de B}}$, con ingreso de
$\boxed{30400\text{ euros}}$.
